\appendix

\stepcounter{appendixnum}
\clearpage
\phantomsection
\chapter*{ПРИЛОЖЕНИЕ А}
\label{appendix:program-guides}
\addcontentsline{toc}{chapter}{ПРИЛОЖЕНИЕ А. Описание программы и руководства}

\begin{center}
\textbf{Описание программы и руководства пользователя и разработчика}
\end{center}

\setcounter{section}{0}
\setcounter{subsection}{0}
\setcounter{subsubsection}{0}
\setcounter{table}{0}
\renewcommand{\thesection}{А.\arabic{section}}
\renewcommand{\thesubsection}{\thesection.\arabic{subsection}}
\renewcommand{\thesubsubsection}{\thesubsection.\arabic{subsubsection}}
\renewcommand{\thetable}{А.\arabic{table}}

\section{Описание программы}

Программа BioSonification предназначена для генерации символической музыки по биологическим последовательностям. Входными данными являются записи в формате FASTA, содержащие DNA, RNA или protein последовательности. Выходными данными являются MIDI-файл с двумя дорожками и, при использовании веб-интерфейса, JSON-метаданные генерации.

Система реализует трехэтапный конвейер:
\begin{enumerate}
\item анализ биологической последовательности и извлечение признаков;
\item генерация гармонической сетки по биологическим признакам;
\item генерация одноголосной мелодии поверх полученной гармонии и запись двухдорожечного MIDI.
\end{enumerate}

Для длинных последовательностей используется фрагментированная генерация: входная последовательность разбивается на участки, каждый участок порождает короткий музыкальный фрагмент, после чего фрагменты объединяются в одну композицию. Такой режим удерживает каждый сгенерированный участок внутри распределения данных, на котором обучалась модель.

Основные функции программы:
\begin{itemize}
\item чтение и очистка FASTA-записей;
\item определение типа последовательности;
\item вычисление биологических признаков и управляющего профиля;
\item сопоставление биологических и музыкальных сегментов при обучении;
\item обучение моделей гармонии и мелодии;
\item генерация MIDI из FASTA через CLI или веб-интерфейс;
\item расчет технических метрик качества сгенерированных MIDI.
\end{itemize}

Минимальные требования для запуска: Python 3.10 или новее, зависимости из файла \texttt{requirements.txt}. Для обучения и быстрой генерации рекомендуется NVIDIA GPU с поддержкой CUDA и не менее 6 GB VRAM. Генерация также может выполняться на CPU, но медленнее.

\section{Руководство пользователя}

Для установки необходимо создать виртуальное окружение и установить зависимости проекта:

\begin{verbatim}
python -m venv .venv
.\.venv\Scripts\python.exe -m pip install --upgrade pip
.\.venv\Scripts\python.exe -m pip install torch --index-url
  https://download.pytorch.org/whl/cu126
.\.venv\Scripts\python.exe -m pip install -r requirements.txt
\end{verbatim}

Генерация музыки из FASTA-файла выполняется командой:

\begin{verbatim}
.\.venv\Scripts\python.exe generate_from_fasta_v2_fragmented.py `
  --checkpoint CHECKPOINT_PATH `
  --fasta data\fasta\refseq_genomes\GCF_000005845.2_genomic.fna `
  --output generated_music.mid `
  --bars-per-fragment 4
\end{verbatim}

Здесь \texttt{CHECKPOINT\_PATH} обозначает путь к файлу \texttt{structured\_pipeline.pt} актуальной обученной модели.

После выполнения команды создается MIDI-файл, который можно открыть в MIDI-плеере, секвенсоре или нотном редакторе. Для примера длиной 10000 bp система создает композицию из нескольких 4-тактовых фрагментов, объединенных в один MIDI.

Для работы через веб-интерфейс используется команда:

\begin{verbatim}
.\.venv\Scripts\python.exe -m web.app
\end{verbatim}

После запуска нужно открыть адрес \texttt{http://localhost:5001}. В веб-интерфейсе можно вставить FASTA-последовательность, загрузить FASTA-файл, запустить генерацию, просмотреть метаданные результата и скачать MIDI. Если настроен аудиосинтезатор, веб-интерфейс также позволяет прослушивать примеры.

Формат входных данных должен соответствовать FASTA: первая непустая строка начинается с символа \texttt{>}, далее следуют строки последовательности. Минимальная длина очищенной последовательности определяется конфигурацией модели.

\section{Руководство разработчика}

Основная реализация находится в пакете \texttt{bio\_music\_pipeline.v2}. Назначение ключевых модулей приведено в таблице~\ref{tab:appendix-modules}.

\begin{table}[htbp]
\centering
\caption{Ключевые модули программы}
\label{tab:appendix-modules}
\begin{tabular}{|p{0.38\textwidth}|p{0.52\textwidth}|}
\hline
\textbf{Модуль} & \textbf{Назначение} \\
\hline
\texttt{bio.py} & Кодирование FASTA, определение типа последовательности, расчет биологических признаков. \\
\hline
\texttt{structured\_music.py} & Извлечение гармонии и мелодии, токенизация, обратный рендеринг в MIDI. \\
\hline
\texttt{structured\_pairing.py} & Сопоставление биологических фрагментов и музыкальных сегментов через дескрипторы. \\
\hline
\texttt{structured\_model.py} & Bio-conditioned autoregressive Transformer для harmony и melody моделей. \\
\hline
\texttt{structured\_train.py} & Оркестрация обучения, сохранение checkpoint, метрик и артефактов. \\
\hline
\texttt{structured\_generate.py} & Инференс: FASTA -> гармония -> мелодия -> MIDI. \\
\hline
\texttt{evaluate.py} & Расчет метрик качества и построение baseline для сравнения. \\
\hline
\end{tabular}
\end{table}

Обучение модели запускается из корня репозитория:

\begin{verbatim}
.\.venv\Scripts\python.exe train_bio_music_v2.py `
  --config configs\pipeline_v2_medium_rtx2060_fast.json `
  --output_dir results\my_experiment
\end{verbatim}

После обучения ожидаются конфигурация запуска, calibration-файл pairing, checkpoint моделей гармонии и мелодии, объединенный \texttt{structured\_pipeline.pt}, JSON-метрики и smoke MIDI. Для проверки корректности изменений рекомендуется запускать тесты и статические проверки:

\begin{verbatim}
.\.venv\Scripts\python.exe -m pytest
.\.venv\Scripts\python.exe -m black . --check
.\.venv\Scripts\python.exe -m flake8 bio_music_pipeline web tests
\end{verbatim}

Веб-интерфейс использует \texttt{web/app.py} и \texttt{web/generator.py}. По умолчанию выбирается checkpoint актуальной модели из каталога \path{results/v2_medium_rtx2060_fast/checkpoints}; при необходимости путь можно переопределить переменной окружения \texttt{BIOSONIFICATION\_STRUCTURED\_CHECKPOINT}.

\section{Дополнительные экспериментальные материалы}

В таблице~\ref{tab:appendix-training-short} приведены основные значения validation loss для актуальной модели. Полные отчеты и промежуточные артефакты сохраняются в каталоге \texttt{results}.

\begin{table}[htbp]
\centering
\caption{Ключевые значения качества обучения}
\label{tab:appendix-training-short}
\begin{tabular}{|p{0.46\textwidth}|c|}
\hline
\textbf{Показатель} & \textbf{Значение} \\
\hline
Harmony validation loss & 0.1454 \\
\hline
Melody validation loss & 0.1566 \\
\hline
Harmony test loss & 0.1389 \\
\hline
Melody test loss & 0.1650 \\
\hline
\end{tabular}
\end{table}

\begin{table}[htbp]
\centering
\caption{Сравнение моделей по chord-tone ratio}
\label{tab:appendix-chord-tone}
\begin{tabular}{|p{0.46\textwidth}|c|}
\hline
\textbf{Вариант} & \textbf{Chord-tone ratio} \\
\hline
Актуальная модель & 0.6710 \\
\hline
Предыдущая версия модели & 0.6073 \\
\hline
Random baseline & 0.4043 \\
\hline
\end{tabular}
\end{table}

\begin{table}[htbp]
\centering
\caption{Пример control profile для фрагмента E. coli}
\label{tab:appendix-control-profile}
\begin{tabular}{|p{0.46\textwidth}|c|}
\hline
\textbf{Параметр} & \textbf{Значение} \\
\hline
Tempo & 0.62 \\
\hline
Density & 0.54 \\
\hline
Harmony change & 0.71 \\
\hline
Register & 0.48 \\
\hline
Complexity & 0.66 \\
\hline
Mode & 0.52 \\
\hline
\end{tabular}
\end{table}
